\begin{table}[H]
\centering
\caption{Modelo ARIMA($p$,0,$q$) del componente estocástico: orden por AIC sobre $p,q\in\{0,\ldots,4\}$, coeficientes principales y diagnóstico de residuos (Ljung--Box).}
\label{tab_arima}
\resizebox{\textwidth}{!}{%
\begin{tabular}{llccccccccc}
\toprule
\textbf{Estación} & \textbf{p} & \textbf{d} & \textbf{q} & \textbf{AR1 (EE)} & \textbf{MA1 (EE)} & \textbf{AIC} & \textbf{BIC} & \textbf{LB p (24)} & \textbf{LB p (168)} & \textbf{Dist. mín. raíces AR-MA} \\
\midrule
ATE & 4 & 0 & 4 & +1.7125 (0.0130) & -0.7835 (0.0145) & 6993.7 & 7053.1 & <0.0001 & <0.0001 & 0.011 \\
\midrule
CDM & 4 & 0 & 4 & +1.6407 (0.0139) & -0.9265 (0.0124) & 7990.3 & 8049.7 & <0.0001 & <0.0001 & 0.007 \\
\midrule
SB & 3 & 0 & 4 & +0.8162 (0.0834) & +0.0472 (0.0835) & 5900.7 & 5953.4 & <0.0001 & <0.0001 & 0.024 \\
\midrule
STA & 4 & 0 & 3 & -0.0134 (0.0150) & +0.7630 (0.0120) & 10387.5 & 10440.3 & <0.0001 & <0.0001 & 0.013 \\
\bottomrule
\end{tabular}}
\par\smallskip{\footnotesize Distancia mínima entre raíces AR y MA < 0.1 indica cancelación aproximada (sobreparametrización).}
\end{table}
